\section{Related Work}
\label{sec:related}

\paragraph{Deep research agents and their evaluation.} Deep research agents extend ReAct \citep{yao2023react} to long-horizon information seeking by interleaving search, document fetching, and synthesis \citep{Xi2025ASurveyof, Shi2025DeepResearchA}; recent systems include proprietary \citep{kimi2026k26} and open variants \citep{Team2025TongyiDeepResearchTechnical, Du2026OpenSeekerDemocratizingFrontier, Li2026OpenResearcherAFully, Li2025WebSailorNavigatingSuper-human} alongside search-enhanced reasoning frameworks \citep{Li2025SearchO1Agentic, Li2025WebThinkerEmpoweringLarge}, with a parallel training literature on trajectory synthesis, RL for search-augmented reasoning, and tool-call efficiency \citep{Jin2025Search-R1TrainingLLMs, Chen2025ReSearchLearningto, Wang2025OTCOptimalTool}. Progress is reported as end-to-end accuracy on BrowseComp \citep{Wei2025BrowseCompASimple}. Our work is orthogonal: rather than building or evaluating an agent end-to-end, we open the trajectory and ask which steps pay off.

\paragraph{Analyzing agent search behavior.} Recent works examine what happens \emph{during} search rather than only at its endpoint. \citet{Ning2026AgenticSearchin} characterize $14$M$+$ requests from DeepResearchGym \citep{Coelho2025DeepResearchGymAFree} at log scale but without per-query relevance supervision; \citet{Chen2025BrowseComp-PlusAMore} provide the human-annotated qrels we build on, but use them to compare retrieval backends rather than diagnose the agent's process; AgentIR \citep{Chen2026AgentIRReasoning-AwareRetrieval} shows reasoning traces are themselves retrieval signals, motivating the snippet-first reading our findings support; other process-level work studies confidence / stopping calibration and trajectory-level evaluation \citep{Ou2025BrowseConfConfidence-GuidedTest-Time, Sun2026DeepSearchwith, Chen2026TRACETrajectory-AwareComprehensive}, and reasoning--action pathologies in which agents over-reason relative to acting \citep{Cuadron2025DangerOverthinking}. None combines the three properties our diagnosis requires --- \emph{per-query document-level qrels}, \emph{multiple agents}, and a \emph{single shared harness and retriever} --- the conjunction that makes the retrieval-vs.-utilization attribution \emph{deterministic} (decided by qrels rather than an LLM judge) and comparable across agents.

% \subsection{Classical IR Analyses of Search Behavior}

% The information-retrieval literature has long studied how humans search, click, and decide to stop, and provides the conceptual scaffolding we adopt for the agent setting. Berrypicking \citep{bates1989berrypicking} models information needs as evolving over the course of a session as new evidence reshapes the question; the query-log analysis of \citet{huang2009analyzing} quantifies how human users reformulate queries when initial results disappoint. Click behavior in human sessions is treated as costly and informative about relevance \citep{joachims2005clickthrough}, motivating our framing of \texttt{visit} as a click-precision metric. Stopping decisions in human search are governed by satisficing and cognitive cost rather than by hard caps \citep{zach2005enough}, in contrast to agents that stop only at context- or step-limit boundaries. Our diagnostic axes --- episode utility, evidence localization, click precision, and the retrieval-vs.-utilization gap --- are direct agent-side analogues of these constructs, applied to trajectories with per-query relevance supervision.
